\section{Lemonade Change}
\label{sec:lemonade-change}

\problemstatement{Customers stand in a queue to buy lemonade, and each glass
costs exactly \$5. We are given an array \textit{bills}, where
$\textit{bills}[i]$ is the bill the $i$-th customer pays with -- always a
\$5, \$10, or \$20 note. We start with no change at all. For every customer,
\textbf{immediately} after they pay, we must hand back the correct change
using bills we are currently holding (we cannot use a bill the same customer
just gave us as their own change). We must process customers strictly in the
given order and decide whether it is possible to give correct change to
\emph{every} customer; return \texttt{true} if so, and \texttt{false} the
moment it becomes impossible.}

\begin{patternbox}
Two phrases give this one away. First, \emph{``customers in a queue,
processed one at a time, in order''} tells us this is an \emph{online
simulation} -- we are never allowed to look ahead or rearrange the input, so
any algorithm that tries to plan globally (sort, search over future choices)
is solving the wrong problem. Second, \emph{``immediately give change''}
combined with only three possible bill values means that at any moment the
entire relevant state of the world is just two small counters: how many \$5
notes and how many \$10 notes we are holding (\$20 notes are never given as
change, so we never need to count them). Whenever a problem reduces to
``maintain a couple of counters and react instantly to each new input,'' it
is a simulation problem with, at most, one local greedy decision per step --
not a search problem.
\end{patternbox}

\subsection*{Understanding the problem carefully}

Since lemonade costs \$5, the change owed depends only on what the customer
pays with:

\begin{itemize}
  \item paying with \$5 $\Rightarrow$ no change owed;
  \item paying with \$10 $\Rightarrow$ owe \$5;
  \item paying with \$20 $\Rightarrow$ owe \$15.
\end{itemize}

A \$5 note can only ever be used to pay back \$5 of change, never \$15
worth by itself, so it has exactly one job: covering a \$10 customer, or
participating in the \$15 change for a \$20 customer. A \$10 note is even
more restricted: it is completely useless for a \$10 customer (we cannot
give a customer back their own denomination as part of correct change
composition here -- effectively, the only practical use of a \$10 note as
change is as part of the \$15 owed to a \$20 customer, paired with one \$5
note). This asymmetry is the heart of the problem: \textbf{\$5 notes are
strictly more flexible than \$10 notes}, because a \$5 note can help pay
back a \$10 customer \emph{or} a \$20 customer, while a \$10 note can only
ever help pay back a \$20 customer.

\subsection*{The only real decision point}

Customers paying with \$5 require no decision -- we simply keep the note.
Customers paying with \$10 also require no decision -- if we have a \$5 note
we must use it (it is the only way to make \$5 change), and if we do not,
failure is immediate and unavoidable, since no other combination of notes we
might be holding sums to exactly \$5 with the unit notes available.

The only point where two different valid strategies exist is a customer
paying with \$20, who is owed \$15. There are exactly two ways to assemble
\$15 from the notes we hold:

\[
  \$10 + \$5 \qquad \text{or} \qquad \$5 + \$5 + \$5.
\]

\begin{ideabox}[Greedy Choice]
At a \$20 customer, if we hold at least one \$10 note and at least one \$5
note, prefer to give \textbf{\$10 + \$5} over \textbf{\$5 + \$5 + \$5}. Only
fall back to spending three \$5 notes if we do not have a spare \$10 note
available.
\end{ideabox}

\subsection*{Why this greedy choice is safe}

Both options return exactly \$15, so the customer is equally satisfied
either way; the difference is entirely in what we are left holding
afterwards. Giving \$10+\$5 costs us one \$5 note. Giving \$5+\$5+\$5 costs
us three \$5 notes. Since \$5 notes are the more broadly useful resource --
the \emph{only} resource capable of covering a future \$10 customer, and a
necessary ingredient (alongside a \$10 note) for covering a future \$20
customer -- depleting three of them at once when we could have depleted only
one is never better for the future, and is sometimes strictly worse: it
might leave us without enough \$5 notes to cover a \$10 customer who arrives
later. Formally, suppose an alternative strategy chooses \$5+\$5+\$5 for some
\$20 customer while a \$10 note sits unused in hand. We can always replace
that choice with \$10+\$5 instead, ending with two extra \$5 notes in hand
compared to the alternative strategy at every subsequent point in the queue,
and strictly more \$5 notes in hand can never cause a future failure that
would not also have occurred with fewer -- so the greedy choice dominates.
\$10 notes, by contrast, have no other use, so spending one whenever
possible is pure upside.

\subsection*{Algorithm}

\begin{enumerate}
  \item Initialize $\textit{five} \gets 0$, $\textit{ten} \gets 0$.
  \item For each bill $b$ in \textit{bills}, in order:
  \begin{enumerate}
    \item If $b = 5$: $\textit{five} \gets \textit{five} + 1$.
    \item If $b = 10$: if $\textit{five} = 0$, return \texttt{false};
          otherwise $\textit{five} \gets \textit{five} - 1$ and
          $\textit{ten} \gets \textit{ten} + 1$.
    \item If $b = 20$:
    \begin{itemize}
      \item if $\textit{ten} \ge 1$ and $\textit{five} \ge 1$: subtract one
            from each;
      \item else if $\textit{five} \ge 3$: subtract three from
            \textit{five};
      \item else: return \texttt{false}.
    \end{itemize}
  \end{enumerate}
  \item If every customer was processed without returning \texttt{false},
        return \texttt{true}.
\end{enumerate}

\begin{algorithm}[H]
\caption{Lemonade Change}
\begin{algorithmic}[1]
\Procedure{LemonadeChange}{\textit{bills}}
  \State $\textit{five} \gets 0$, $\textit{ten} \gets 0$
  \For{each $b$ in \textit{bills}}
    \If{$b = 5$}
      \State $\textit{five} \gets \textit{five} + 1$
    \ElsIf{$b = 10$}
      \If{$\textit{five} = 0$} \Return \textbf{false} \EndIf
      \State $\textit{five} \gets \textit{five} - 1$
      \State $\textit{ten} \gets \textit{ten} + 1$
    \Else \Comment{$b = 20$}
      \If{$\textit{ten} \ge 1$ \textbf{and} $\textit{five} \ge 1$}
        \State $\textit{ten} \gets \textit{ten} - 1$
        \State $\textit{five} \gets \textit{five} - 1$
      \ElsIf{$\textit{five} \ge 3$}
        \State $\textit{five} \gets \textit{five} - 3$
      \Else
        \State \Return \textbf{false}
      \EndIf
    \EndIf
  \EndFor
  \State \Return \textbf{true}
\EndProcedure
\end{algorithmic}
\end{algorithm}

\subsection*{Dry run}

\dryrun{Take $\textit{bills} = [5, 5, 5, 10, 20]$.}

Start with $\textit{five}=0$, $\textit{ten}=0$.

\begin{itemize}
  \item Bill $5$: $\textit{five}=1$.
  \item Bill $5$: $\textit{five}=2$.
  \item Bill $5$: $\textit{five}=3$.
  \item Bill $10$: $\textit{five} > 0$, so give one \$5. $\textit{five}=2$,
        $\textit{ten}=1$.
  \item Bill $20$: $\textit{ten}\ge1$ and $\textit{five}\ge1$, so give
        \$10+\$5. $\textit{five}=1$, $\textit{ten}=0$.
\end{itemize}

All five customers are served, so the answer is \texttt{true}.

Now take $\textit{bills} = [5, 5, 10, 10, 20]$.

\begin{itemize}
  \item Bill $5$: $\textit{five}=1$.
  \item Bill $5$: $\textit{five}=2$.
  \item Bill $10$: give one \$5. $\textit{five}=1$, $\textit{ten}=1$.
  \item Bill $10$: give one \$5. $\textit{five}=0$, $\textit{ten}=2$.
  \item Bill $20$: need $\textit{ten}\ge1$ and $\textit{five}\ge1$, but
        $\textit{five}=0$. Fall back to $\textit{five}\ge3$, but again
        $\textit{five}=0$. Neither option works, so return \texttt{false}.
\end{itemize}

\subsection*{C++ implementation}

\begin{lstlisting}
#include <bits/stdc++.h>
using namespace std;

bool lemonadeChange(vector<int>& bills) {
    int five = 0, ten = 0;

    for (int bill : bills) {
        if (bill == 5) {
            five++;
        } else if (bill == 10) {
            if (five == 0) return false;
            five--;
            ten++;
        } else { // bill == 20
            if (ten >= 1 && five >= 1) {
                ten--;
                five--;
            } else if (five >= 3) {
                five -= 3;
            } else {
                return false;
            }
        }
    }
    return true;
}

int main() {
    vector<int> bills = {5, 5, 5, 10, 20};
    cout << (lemonadeChange(bills) ? "true" : "false") << endl;
    return 0;
}
\end{lstlisting}

\begin{complexitybox}
\textbf{Time Complexity.} Each of the $n$ customers is processed exactly
once, with only constant work (a couple of comparisons and increments) per
customer, giving
\[
  O(n).
\]
\textbf{Space Complexity.} We track only the two counters \textit{five} and
\textit{ten}, so the space complexity is $O(1)$.
\end{complexitybox}

\begin{takeawaybox}
When several local choices return the same immediate result (here, both
change-giving options return exactly \$15), break the tie by asking which
choice preserves the resource that is more broadly useful for the future.
Here, \$5 notes serve two future needs while \$10 notes serve only one, so
spending a \$10 note whenever possible is always at least as good, and
sometimes strictly necessary, for surviving the rest of the queue.
\end{takeawaybox}
